\section*{Limitations}
\label{sec:limitations}
\addcontentsline{toc}{section}{Limitations}

\paragraph{The environment is engineered to create the trade-off.} The counterparty shows its best overall offer early and by construction that offer breaks the client's requirement, which is why a blocked better deal exists in \exposureRate\ of episodes. This is a design choice that makes the measurement possible; it is not an estimate of how often such a trade-off arises in real delegated work. The rates here answer ``when a blocked better deal exists, how often is it conveyed'', not ``how often does an agent have something like this to convey''. The environment is also one class of bilateral procurement with an explicit point table, which makes value errors checkable but also makes arithmetic part of the task, and the counterparty is scripted---realistic in structure, not in language.

\paragraph{Behaviour, not intent.} The design cannot separate concealment from failure to track, and we do not use ``deception'' for any finding. A correct knowledge-probe answer shows the information was recoverable on request, not that it was salient while the report was being written. The preregistration commits to this reading and we hold to it.

\paragraph{Sampling and reproducibility.} Each cell is run once at provider-default decoding, so within-cell sampling variance is not separately estimated: every interval here covers variation across scenarios, not across repeated samples of one episode. Models are served behind aliases that providers may repoint; the served-model strings recorded per episode, not the product names, are the reproducible identifiers, and collection ran inside a two-day window to limit drift.

\paragraph{Validation is thinner than planned.} The preregistration specified 150 disclosure, 150 characterization, 100 stated-total and 50 leak items. Two annotators labelled 50 disclosure and 25 stated-total items; characterization ratings and leak detection received no human validation at all and are reported as secondary throughout. Because reports that convey the blocked deal are rare, the validation set contains few of them, so the judge's false-positive rate is weakly identified; \cref{sec:validation} reports a sensitivity analysis across the range the labels leave open rather than a point estimate. One of the two annotators is an author.

\paragraph{The principal is simulated.} H5 asks what a principal decides from a report, and the principal is a language model that received no validation against human readers. It is a claim about how one model responds to these reports, not about how a person would. Whether human principals would sign off differently is untested and is the natural next study.

\paragraph{Measurement floors.} The judge-free measures are deliberately broad and therefore lower bounds: a report can name the blocked offer without conveying what it was worth. Leak detection is verbal only, so a confidential figure could in principle leak through the pattern of concessions without being counted. The stated-total extractor misses about one claim in six, so the rate at which models state a value is itself a floor.
